\section{Design}
\label{sec:design}

\subsection{Overview}

VGuide intercepts the CEGAR loop at a single well-defined point: immediately after the
first spurious counterexample is confirmed infeasible, before standard interpolation
runs. The intervention is deliberate: the first spurious CE is the earliest moment at
which the CEGAR engine has identified a concrete execution path that fails to satisfy
the property, giving the LLM a rich, program-specific context to reason about. At this
point, interpolation-based refinement would add only one locally-scoped predicate.
VGuide instead builds a structured \emph{context pack}, queries the LLM for candidate
predicates, filters the response through a three-tier validation cascade, and injects
all surviving predicates simultaneously at CFA loop heads. Standard CEGAR refinement
then resumes and handles all subsequent spurious CEs through interpolation as usual.
The pipeline is illustrated in \Cref{fig:pipeline}.

\begin{figure}[t]
  \centering
  \begin{tikzpicture}[
    font=\small,
    box/.style={draw, rounded corners=3pt, minimum width=2.0cm, minimum height=0.75cm,
                align=center, inner sep=4pt},
    vbox/.style={draw, rounded corners=3pt, minimum width=2.2cm, minimum height=0.65cm,
                 align=center, inner sep=3pt, fill=gray!12},
    region/.style={draw=gray!60, fill=gray!8, rounded corners=5pt, dashed},
    arr/.style={->, >=stealth, thick},
    annot/.style={font=\scriptsize, text=gray!70},
  ]

  \node[box] (prog) at (0,0) {Program\\(C)};

  \node[box, minimum width=2.2cm, minimum height=1.6cm] (cegar) at (3.2,0) {};
  \node[font=\small\bfseries] at (3.2,0.2) {CEGAR};
  \node[annot] at (3.2,-0.45) {(CPAChecker)};

  \node[box, fill=blue!8, minimum width=2.2cm] (bridge) at (6.6,0) {VGuide\\Bridge};

  \node[vbox] (ctx)  at (10.2, 1.0) {Context Pack};
  \node[vbox] (llm)  at (10.2, 0.0) {LLM\\(\LLMName{})};
  \node[vbox] (val)  at (10.2,-1.0) {L1 / L2 / L3\\Validation};

  \begin{scope}[on background layer]
    \node[region, fit=(ctx)(llm)(val), inner sep=6pt] (vregion) {};
  \end{scope}

  \draw[arr] (prog.east) -- (cegar.west)
    node[midway, above, annot] {};

  \draw[arr] ([yshift=3pt]cegar.east) -- ([yshift=3pt]bridge.west)
    node[midway, above, annot] {Spurious CE};

  \draw[arr] ([yshift=-3pt]bridge.west) -- ([yshift=-3pt]cegar.east)
    node[midway, below, annot] {Valid preds};

  \draw[arr] (bridge.east) -- (ctx.west);

  \draw[arr] (ctx.south) -- (llm.north);
  \draw[arr] (llm.south) -- (val.north);

  \draw[arr] (val.west) to[out=180, in=0] (bridge.east);

\end{tikzpicture}

  \caption{VGuide pipeline. CEGAR fires VGuide on the first spurious counterexample.
           Valid predicates are injected at loop heads; all others are silently
           discarded.}
  \label{fig:pipeline}
\end{figure}

This design has two properties. First, it is \emph{conservative}: if VGuide injects
no valid predicates---because the LLM response is empty, malformed, or entirely
filtered---the CEGAR loop continues exactly as stock, with no behavior change. Second,
it is \emph{sound}: every predicate that reaches the abstract domain has been certified
to be a syntactically valid C Boolean expression over variables in scope, consistent
with the current abstract state. The soundness guarantee of CEGAR is preserved
unconditionally.

\subsection{Context Pack}

The context pack is the structured input to the LLM and contains three components.

\paragraph{Stripped source.} The C source of the verification task is preprocessed to
remove boilerplate: \texttt{\#include} directives, non-verification macros, and
commented-out code are stripped. The remaining source is formatted to fit within the
LLM's context window. This gives the model access to the global program structure
including variable declarations, function signatures, and loop topology.

\paragraph{CE trace in block formula format.} The spurious CE is translated into
SMT-LIB2 block formula format, a sequential conjunction of path formulas
$\varphi_1 \wedge \varphi_2 \wedge \cdots \wedge \varphi_k = \bot$ where each
$\varphi_i$ represents one basic-block transition. SSA renaming is applied so that
all variable versions are distinct. The block formula makes the quantitative loop
state at each CE step explicit: the LLM can read off the values of \texttt{x} and
\texttt{y} at each iteration and observe the relational pattern.

\paragraph{CFA loop-head labels.} The control-flow automaton (CFA) node identifiers
for all loop heads in the program are extracted and serialized as a JSON object whose
keys are the loop-head node names. This serves a dual purpose: it provides the LLM
with explicit injection targets for its predicate suggestions, and it constrains the
output format. The LLM is required to return predicates keyed by these exact node
names; any response that references a node name not in the JSON key set is rejected
at L1 validation. Hallucinated node names cannot survive.

\paragraph{Dual-prompt design.} VGuide uses two prompt variants that are sent
independently and whose results are merged before validation. The \emph{SAFE variant}
instructs the LLM to produce loop invariants that, if true, support a safety proof
for the property under analysis. The \emph{BUG variant} instructs the LLM to produce
predicates that distinguish safe from unsafe execution paths, useful for
over-approximating the reachable error region. Merging both prompt responses maximizes
the diversity of candidate predicates submitted to the validation cascade.

\subsection{LLM Choice and Configuration}

VGuide uses \LLMName{} as its predicate oracle. The choice is motivated by three
factors. First, DeepSeek V4 Pro achieves 80.6 on the SWE-bench Verified benchmark,
demonstrating strong performance on structured code reasoning tasks. Second, at
\$0.87 per million output tokens it is cost-effective for use across hundreds of
verification tasks. Third, and most importantly for VGuide's operating context,
\LLMName{} supports disabling its chain-of-thought \emph{thinking} mode. Standard
thinking mode produces long reasoning traces before the answer, incurring 52--88\,s
of latency per call. Since predicate expressions are short C Boolean formulas that
require no multi-step derivation visible to the caller, thinking adds latency without
adding usable output. With thinking disabled, \LLMName{} returns a JSON predicate
map with a median latency of \LLMLatencyMedian{}\,s. Within a 300\,s per-task
budget, this makes multiple scheduling rounds feasible and ensures that the LLM
overhead is a minor fraction of the total analysis time.

\subsection{Three-Tier Validation Cascade}

All LLM-proposed predicates pass through three validation tiers before reaching the
CEGAR abstraction. A predicate is accepted only if it clears all three tiers; failure
at any tier results in silent discard.

\paragraph{L1: Contract validation.} The \texttt{PredicateContractValidator} checks
the structural contract of the LLM's response. It verifies that the response is valid
JSON, that each top-level key is a recognized CFA loop-head node name (from the list
supplied in the context pack), and that each value in the associated predicate array is
a non-empty string. It further checks that each predicate string parses as a syntactically
valid C Boolean expression using a lightweight grammar checker. Predicates containing
assignment operators (\texttt{=} without equality semantics), side-effecting function
calls, or dereferencing operators on non-pointer types are rejected. This tier filters
hallucinated node names, JSON schema violations, and syntactically malformed expressions.

\paragraph{L2: Vocabulary filtering.} The \texttt{VocabularyGuide} maps each loop
head to the set of program variables that are in scope at that location, as determined
by the CFA. A predicate proposed for a given loop head is accepted only if every free
variable appearing in the predicate belongs to the loop head's vocabulary set. Predicates
referencing variables that do not exist at the injection point---a common LLM
hallucination pattern---are silently dropped. This tier prevents type errors and
undeclared-variable errors that would otherwise surface as CPAChecker internal
exceptions.

\paragraph{L3: SMT entailment.} For each predicate that passes L1 and L2, VGuide
invokes Z3 to check two conditions. First, the predicate is checked for satisfiability:
a predicate that is trivially false (e.g., $x > 0 \wedge x < 0$) is discarded as it
would collapse the abstract state at the injection point. Second, the predicate is
checked for non-redundancy: if the predicate is already entailed by the current abstract
state at the loop head, adding it to $\Pi$ would not change the abstraction, and it is
discarded to avoid polluting the predicate set with no-ops. Only predicates that are
satisfiable and not already entailed are accepted. This tier ensures that every accepted
predicate represents a genuine refinement of the abstract domain.

\subsection{Injection and Soundness}

Predicates that pass L1, L2, and L3 are added to the precision of the PredicateCPA
at the corresponding CFA loop-head locations. CPAChecker's precision update mechanism
triggers re-abstraction at those locations, and the CEGAR loop continues from the
refined abstract state.

The soundness argument is as follows. The PredicateCPA maintains the invariant that
the abstract domain is always a sound over-approximation of the concrete program
states: for any concrete state $s$ reachable in $P$, the corresponding abstract state
$\hat{s}$ satisfies $\hat{s} \supseteq \alpha(s)$ where $\alpha$ is the abstraction
function. Adding new predicates to the precision strictly refines the abstraction:
more predicates mean a finer abstract domain, and the over-approximation invariant is
preserved. The critical point is that VGuide adds predicates only; it never removes
predicates or modifies the model-checking or CE-feasibility-checking components. The
only risk of unsoundness would be if a predicate injection caused the abstract domain
to no longer be a sound over-approximation---but this is ruled out because (i) L3
ensures all predicates are satisfiable, and (ii) adding more predicates to a predicate
abstraction can only make it more precise, never less. VGuide therefore cannot introduce
a wrong-verdict without a pre-existing soundness bug in CPAChecker's PredicateCPA.

In the VGuide architecture, the LLM is classified as a Tier~S (suggestion) component:
it provides hints to the verification engine but has no authority over the abstract
state, the refinement decision, or the final verdict. The CEGAR engine remains the
sole decision maker.
